\clearpage
\section{When do deeper peptide models help against additive baselines?}
\subsection{The two estimands must stay separate}
The peptide-disjoint IEDB test asks whether a model ranks held-out peptides for alleles present in the training context. The locked G3 challenge asks a different question: ranking for held-out alleles. A good same-row score on familiar alleles cannot be used to answer the second question. On the peptide-disjoint test, the per-allele PSSM AUROC is 0.9156, the allele-conditioned CNN 0.9044 and the GNN 0.8480. A z-scored PSSM/CNN ensemble attains 0.9273 on that familiar-allele split; this comparison does not establish a deep-model improvement over the additive model alone on its primary row set. The CPP task also favors a simpler baseline in the recorded test: the 3-mer logistic AUROC is 0.9265 versus 0.8956 for the CNN. These are two different biological prediction tasks, not pooled evidence for a universal representation advantage.

The selected 6,000-row familiar-allele head-to-head reports ensemble AUROC 0.9281 versus MHCflurry 0.9164. This is a local paired comparison, not a new-allele transfer certificate. Its high binder prevalence affects the interpretation of PPV; low-FPR behavior, calibration and errors on allele shift remain separate questions. Thirty allele-specific comparisons within this panel are correlated because alleles share assay design and peptide context; a nominal per-allele Wilcoxon result cannot turn them into independent cohorts.

\subsection{The locked G3 result overrides the broad claim}
In \texttt{results/g3\_allele\_holdout.json}, 83,333 training examples and 6,000 evaluations across twelve held-out alleles yield pooled AUROC 0.6687 for the tested model versus 0.9285 for MHCflurry. Mean per-allele AUROC is 0.5798 versus 0.9056. The gate is recorded false. A small 40-row HLA-C*06:02 stratum is individually unstable, but it cannot account for the aggregate gap; all twelve reported allele comparisons numerically favor the comparator. The manuscript therefore does not say that this method beats a leader on unseen alleles. A second independent held-out-allele test remains [PENDING], not an inferred replication of G3.

\subsection{Mechanistic explanations are testable hypotheses}
A per-allele PSSM can exploit assay examples for each familiar allele. An allele-conditioned CNN may learn patterns useful inside the observed allele panel yet lack support for pocket configurations of held-out alleles. The G3 split may also shift peptide length, assay method, binder prevalence and binding motifs. The present aggregate and per-allele score tables alone cannot allocate the failure among these causes. The committed pocket-entropy and groove-geometry audits show diversity in the observed HLA panel, but they are not a controlled intervention on sequence encoding and cannot prove that an unseen pocket geometry caused any specific error. A causal explanation would require predeclared allele-distance and pocket-site strata, train/test assay and peptide overlap audits, uncertainty by allele, and a paired comparison that changes only the representation. Until those tests exist, ``allele shift'' names the evaluation condition rather than a diagnosed molecular mechanism.

\subsection{Decision path and prospective repair}
For a peptide from an allele represented in training, report the per-allele PSSM beside the CNN and ensemble on the same assay and split. Require low-FPR and calibration evidence before using a threshold or top-ranked list. For an allele unseen at training, this tested model should abstain from a reliable binder ranking claim: the G3 gap is too large to import the familiar-allele score. A future attempt should lock the set of unseen alleles, endpoint, comparator versions, inclusion rules and thresholds before scoring, then test stronger allele representations such as protein embeddings on the same rows. A ligand-eluted benchmark answers presentation rather than binding-affinity measurement and must be labeled a new task; it is [PENDING]. An observed-design rank audit is also [PENDING], so the algebraic identifiability condition in the earlier section remains conditional. Counts of assays, peptide-allele rows, alleles, downloaded files and independent datasets must not be collapsed into one number. This section adds interpretation and falsifiable tests, not a new numeric outcome or a slide deliverable.
